
\paragraph{Benchmarking Autonomous Agents}

Table \ref{tab:benchmark_comparison} offers a comparison between \method and existing benchmarks, highlighting the unique contributions of our proposed benchmarking suite. Existing benchmarks for autonomous agents, such as those introduced by \citet{brockman2016openai, bellemare2013arcade, tassa2018deepmind}, provide diverse environments for testing various algorithms. However, these benchmarks often focus on specific types of learning algorithms or on evaluating particular desirable qualities in autonomous agents. For example, \citet{fu2020d4rl} and \citet{rl_unplugged} evaluate offline reinforcement learning \citep{levine2020offline}, while \citet{meta_world} focuses on meta-reinforcement learning \citep{wang2016learning_to_reinforcement_learn}. Similarly, \citet{ye2021mastering} tests sample efficiency, \citet{guss2019minerl} challenges agents on long-horizon tasks, and \citet{coin_run} evaluates generalization ability.
While these benchmarks provide insights, they do not fully capture the challenges faced by autonomous agents in real-world applications \citep{dulac2021challenges}.
Environments, benchmarks, and datasets have been made to foster the development of autonomous agents in real-world scenarios, such as aerial balloon navigation~\citep{balloon_learning_env}, autonomous driving~\citep{waymo_open}, website navigation \citep{gur2021environment}, and furniture assembly \citep{lee2021ikea}. Yet, these initiatives are often domain-specific and lack the comprehensive scope needed to evaluate agents across a wide range of real-world challenges  as outlined by prior work \citep{dulac2021challenges}, which forms the basis for our work.


Consequently, there remains a need for more benchmarking suites that encompass a diverse set of tasks and environments, reflecting the complexity and variety of problems encountered in real-world applications. Among recent benchmarks, NeoRL \citep{qin2022neorl} provides realistic environments for stock trading, utility management and industrial control, while OGBench \citep{park2024ogbench} emphasizes realistic tasks in the offline, goal-conditioned setting. Jumanji \citep{bonnet2023jumanji} focuses on providing fast, JAX-accelerated \citep{jax2018github} implementations of combinatorial optimization tasks inspired by industry applications.

\method differentiates itself by incorporating different real-world domains such as web navigation and computer chip floorplanning, while also including system performance, data cost, and reliability metrics in a unified package. This comprehensive approach allows for a more holistic evaluation of autonomous agents across diverse, practically relevant tasks and crucial deployment considerations.

\paragraph{Benchmarking System Performance} In addition to evaluating task-specific performance metrics, analyzing the end-to-end performance cost and examining the hardware resources required to apply learning algorithms on specific environments has gained significant attention~\citep{wu2022sustainable,Howwe’re47:online}. Benchmarks such as MLPerf \citep{reddi2020mlperf} and DAWNBench \citep{Coleman2017DAWNBenchA} have been developed to assess various aspects of commercial deep learning workloads across training and inference, considering a diverse class of systems. Furthermore, recent studies have investigated the environmental impact of deep learning by quantifying the carbon footprint associated with training and inference using large neural network models \citep{patterson2021carbon}. This line of research has also extended to autonomous agents, with works like QuaRL demonstrating reduced energy consumption and emissions through lower-precision distributed training \citep{krishnan2022quarl}. 
Despite these efforts, there remains a need for evaluating the system performance and energy consumption of autonomous agents to provide valuable insights into their practical feasibility and sustainability.


\paragraph{Reliability Metrics for Reinforcement Learning}
Reliability is a concern in reinforcement learning (RL), as current metrics often rely on point estimates of aggregate performance, which fail to capture the true performance of algorithms and make it challenging to draw conclusions about the state-of-the-art \citep{agarwal2021deep_rl_stat_precipice, henderson2018deep_rl_that_matters, colas2018many}.
The increasing complexity of benchmarking tasks has made it infeasible to run hundreds of training runs, necessitating the development of tools to evaluate reliability based on a limited number of runs \citep{agarwal2021deep_rl_stat_precipice}.
For real-world deployments, reliability is essential to ensure that RL algorithms perform consistently and robustly across different conditions and environments.
To assess reliability, it is essential to consider metrics across three axes of variability: time (within a training run), runs (across random seeds), and rollouts of a fixed policy \citep{chan2019measuring}. By incorporating reliability metrics into \method, we will be able to better assess the robustness and consistency of RL algorithms.
